\section{The Third Week: Rejoice, the Lord Is Near}\label{sec:week-3}

On the Third Sunday of Advent, the Roman Church assembles at the
Basilica of St.~Peter (\textit{Statio ad S.\ Petrum}), and the Mass
opens with a sudden burst of gladness. The Introit takes its text from
St.~Paul's letter to the Philippians: \textit{Gaudete in Domino semper:
iterum dico gaudete}, ``Rejoice in the Lord always: again, I say,
rejoice. Let your modesty be known to all men. The Lord is nigh''
(Phil~4:4--5), paired with the verse from Psalm~84~[85]:2, ``Lord, thou
hast blessed thy land: thou hast turned away the captivity of Jacob.''
From that opening word, \textit{Gaudete}, the Sunday takes its
traditional name.

The 1962 rubrics mark this midpoint of Advent with measured festivity.
Rose vestments are permitted on the Third Sunday of Advent---and only in
the Office and Mass of the Sunday itself, not on the weekdays that
repeat its Mass (RG~n.~131). The organ or harmonium is permitted to
sound on Gaudete Sunday (Sacred Congregation of Rites, \textit{Instructio
de musica sacra}, 3~September 1958, n.~83~b), lifting the musical
restraint of the first two weeks. Yet \textit{Gloria in excelsis} and
\textit{Te Deum} remain silent; Christmas is close, not yet arrived.
Even amid the rose vestments, the chants preserve the longing of
Advent. In the Gradual (\textit{Qui sedes, Domine}, Ps~79~[80]:2--3),
the choir prays, ``Stir up thy might, and come to save us,'' and repeats
the plea in the Sunday Alleluia (\textit{Excita, Domine, potentiam
tuam}, omitted on the ordinary weekdays of the week). At the Offertory
(\textit{Benedixisti, Domine}, Ps~84~[85]:2), the Church recalls that
God has turned away the captivity of Jacob, and at Communion
(\textit{Dicite: pusillanimes}, Is~35:4) she addresses the weary heart:
``Say to the fainthearted: Take courage, and fear not: behold your God
will bring the revenge of recompense: God himself will come and will
save you.''

Where the First, Second, and Fourth Sundays begin their Collects with
\textit{Excita} (``Stir up''), Gaudete Sunday lowers its voice into a
quiet, intimate petition addressed directly to Christ:

\begin{massbox}{Collect of the Third Sunday of Advent (\textit{Gaudete})}
\textbf{Latin.} Aurem tuam, quaesumus, Domine, precibus nostris
accommoda: et mentis nostrae tenebras gratia tuae visitationis illustra:
Qui vivis et regnas cum Deo Patre in unitate.

\textbf{English.} Incline Thine ear to our prayers, O Lord, we beseech
Thee; and make bright the darkness of our minds by the grace of Thy
visitation. Who livest.
\end{massbox}

Here the Church asks Christ (\textit{Qui vivis}, ``Who livest'') to bend
close to our prayers and scatter the darkness of our minds simply by
coming to visit us. The Secret (\textit{Devotionis nostrae}) and
Postcommunion (\textit{Imploramus, Domine}) follow the same line, asking
that the continual offering of this sacrifice may work our salvation
within us and cleanse us from our sins for the coming feast.

\subsection{Joy Born of Contrition: St.~Paul and St.~John Chrysostom}

In the Epistle (Phil~4:4--7), the Church hears again the passage that
opened the Mass, now carried to its close: ``Rejoice in the Lord always:
again, I say, rejoice. Let your modesty be known to all men. The Lord is
nigh. Be nothing solicitous: but in every thing, by prayer and
supplication, with thanksgiving, let your petitions be made known to
God. And the peace of God, which surpasseth all understanding, keep your
hearts and minds in Christ Jesus.''

St.~Paul wrote those words in chains to a community enduring affliction
for the faith. Preaching on this Epistle (\textit{Homily} 14 on
Philippians), St.~John Chrysostom immediately raises the objection any
serious hearer feels: Christ himself pronounced ``Blessed are they that
mourn'' and ``woe unto them that laugh.'' How then can Paul command us
to rejoice always? Chrysostom draws a sharp line between the laughter of
the world and the joy of the Christian: ``This joy is not contrary to
that grief, but from that grief it too is born. For he who grieveth for
his own faults, and confesseth them, rejoiceth. Moreover, it is possible
to grieve for our own sins, and yet to rejoice in Christ.'' Because such
joy rests on communion with God rather than on comfortable
circumstances, neither hardship nor chains can extinguish it.

Chrysostom then examines Paul's twin remedies against anxiety. The first
is the nearness of Christ---``The Lord is nigh''---which steadies the
faithful under injustice because the judgment is already close at hand.
The second is prayer accompanied by thanksgiving, which Chrysostom calls
``a medicine which healeth grief, and distress, and all that is
painful.'' A Christian gives thanks even in trouble, trusting that God
orders all things for our good even when we cannot yet trace how. And
when Chrysostom reaches Paul's closing promise of the peace of God that
surpasses all understanding, he points straight to the Incarnation and
Redemption:

\begin{quote}
For who could have expected, who could have hoped, that such good
things would have come? They exceed all man's understanding, not his
speech alone. For His enemies, for those who hated Him, for those who
determined to turn themselves away, for these, he refused not to deliver
up His Only Begotten Son, that He might make peace with us. This peace
then, i.e.\ the reconciliation, the love of God, shall guard your hearts
and your thoughts.
\end{quote}

\subsection{The Voice and the Word}

In the Gospel (John~1:19--28), priests and Levites come out from
Jerusalem to question John at the Jordan: ``Who art thou?'' St.~John
the Evangelist slows the narrative down with a solemn triple emphasis:
``And he confessed and did not deny: and he confessed: I am not the
Christ.'' Asked whether he is Elias, he says, ``I am not.'' Asked
whether he is the prophet, he answers, ``No.'' Pressed to give an
account of himself to those who sent them, the Baptist narrows his whole
identity to one line from Isaiah~40:3: ``I am the voice of one crying in
the wilderness, make straight the way of the Lord, as said the prophet
Isaias.'' Then he speaks the sentence that searches every Advent
congregation: ``I baptize with water: but there hath stood one in the
midst of you, whom you know not. The same is he that shall come after
me, who is preferred before me: the latchet of whose shoe I am not
worthy to loose'' (John~1:26--27).

Pope St.~Gregory the Great preached his seventh homily on the Gospels
on this passage, and the Roman Breviary reads its opening at Matins on
the Third Sunday of Advent. Gregory stops first over John's refusal to
borrow a dignity that did not belong to him: ``In that he arrogated not
to himself the name of Christ, he became a member of Christ. While he
humbly strove to confess his own weakness, he earned by his simplicity a
part in the grandeur of his Master.'' Then, in the second part of the
same homily (\textit{Homily} 7.2, in the \textit{Catena Aurea} on
John~1:23), Gregory unveils why the Baptist chooses the exact word
\textit{voice} to describe himself:

\begin{quote}
Ye know that the only-begotten Son is called the Word of the Father.
Now we know, in the case of our own utterance, the voice first sounds,
and then the word is heard. Thus John declares himself to be the voice,
i.e.\ because he precedes the Word, and, through his ministry, the Word
of the Father is heard by man.
\end{quote}

Once the Word enters the listener's heart, the passing sound has done
its work. St.~Augustine reads the same interrogation
(\textit{Tractates on John} IV.1, IV.7) with his gaze fixed on the
unrecognized Lord already standing among the crowds. Where earlier prophets foretold Christ from
afar, John was granted ``to point Him out with the finger,'' yet the
crowds could not see him because ``He had come humbly and concealed from
the first; the more concealed in proportion as He was more humble.''
Then Augustine contrasts John with every earthly herald who marches
ahead of a magistrate (\textit{Tractates on John} IV.7):

\begin{quote}
Would it not have seemed to you that a herald would have cried, ``Go
away, make room.'' Instead of the herald's cry ``Go away,'' John says
``Come.'' The herald makes men stand back from the judge; to the Judge
John calls. Yes, indeed, John calls men to the lowly One, that they may
not experience what He will be as the exalted Judge.
\end{quote}

\subsection{At the Threshold of the Ember Days and the Great Antiphons}

The Third Week of Advent stands at the hinge of the season. On Sunday
morning, Matins opens Isaiah~26: ``In that day shall this canticle be
sung in the land of Juda. Sion the city of our strength a saviour, a
wall and a bulwark shall be set therein'' (Is~26:1). Monday reads
Isaiah~28, Tuesday opens Isaiah~30 with the tender promise, ``Therefore
the Lord waiteth that he may have mercy on you: and therefore shall he
be exalted sparing you: because the Lord is the God of judgment: blessed
are all they that wait for him'' (Is~30:18), and Thursday reads
Isaiah~33. On Monday, Tuesday, and Thursday, whenever no feast such as
St.~Lucy (13~December) or St.~Eusebius (16~December) intervenes, the
priest repeats the Mass of Gaudete Sunday---now in violet vestments,
never rose (RG~n.~131), and omitting the \textit{Alleluia} (RGMR~n.~299).
When 13~December falls on Gaudete Sunday itself (in years when Advent
begins on 29~November), the Sunday takes precedence as a first-class day
and St.~Lucy is omitted.

On Wednesday, Friday, and Saturday of this week arrive the Winter Ember
Days---three second-class weekdays with proper Masses, Gospel homilies
at Matins instead of Isaiah (RGBR~n.~222~a), and an ancient Roman
discipline of fasting and almsgiving. And in the shortest Advent, when
the First Sunday falls on 3~December, Gaudete Sunday itself lands on
17~December, so that at Second Vespers of Gaudete Sunday the choir
already chants \textit{O Sapientia}, the first of the Seven Great
Antiphons. In every year, Ember Saturday falls between 17~December and
23~December, inside the week of the Great Antiphons.

In the older nine-lesson Roman Breviary, the Church prepared the
faithful for the coming Ember fast on Gaudete Sunday itself by reading
Pope St.~Leo the Great's second sermon on the December fast
(\textit{Sermon} 13) in the second nocturn of Matins. Though the 1960
Sunday Office limits Matins to Isaiah and St.~Gregory's Gospel homily
(RGBR~n.~220), St.~Leo's words remain the best preparation for the week
ahead:

\begin{quote}
Fasting hath ever been the bread of strength. From abstinence proceed
pure thoughts, reasonable desires, and healthy counsels. By voluntary
mortifications the flesh dieth to lust, and the soul is renewed in
might. But since fasting is not the only mean whereby we get health for
our souls, let us add to our fasting works of mercy. Let us spend in
good deeds what we take from indulgence. Let our fast become the banquet
of the poor.
\end{quote}

Across the opening days of this Third Week, read Philippians~4:4--7,
John~1:19--28, and the appointed chapters of Isaiah (see
Appendix~\ref{app:daily-course} for each day's readings and prayer),
setting aside time with St.~John Chrysostom on Philippians,
St.~Gregory's seventh homily, and St.~Leo's December sermons. Pray the
\textit{Aurem tuam} Collect each day (Appendix~\ref{app:prayers}),
begging Christ to brighten whatever remains dark or anxious in your
mind. Let the joy of Gaudete Sunday take concrete form without clamour:
practice St.~Paul's ``modesty''---forbearance and gentleness toward
everyone you meet---accept small slights without defending your own
reputation as John refused borrowed titles, and decide on Monday or
Tuesday what concrete alms you will set aside for the poor as the Ember
Days begin (see Section~\ref{sec:penance-law} for the discipline of
penance and the Ember week).
